
\section{From heterogeneous inputs to embedded spatial graphs}
\label{supp:representation}

\begin{figure}[!t]
\centering
\includegraphics[width=\linewidth,keepaspectratio]{GeneralVersion/Figures/InputYamadaExamples.pdf}
\caption{\textbf{A common spatial-graph interface supports invariant calculations across different source representations.}
Representative inputs are \textbf{(a)} a trefoil polymer \cite{kg_deguchi2017topological_polymers}; \textbf{(b)} a Petersen spatial network; \textbf{(c)} a genus-two generating graph illustrating the regular-neighborhood viewpoint \cite{kg_ishii2008handlebody,kg_ishii2012handlebody}; \textbf{(d)} M\"obius ladder $M_8$; \textbf{(e)} a figure-eight knot; \textbf{(f)} an engineering network \cite{kg_shai1999engineering_graphs}; \textbf{(g)} a cinquefoil knot; \textbf{(h)} a warped cubical cage; \textbf{(i)} a Lissajous loop; \textbf{(j)} a triple-lens necklace; \textbf{(k)} a spatial $K_{3,3}$ graph; \textbf{(l)} a twisted triangular prism; \textbf{(m)} a three-component $T(3,3)$ torus link; \textbf{(n)} a spatial network originally stored in GraphML \cite{kg_brandes2002graphml} and converted externally; and \textbf{(o)} a truss network \cite{kg_shai1999engineering_graphs}. Red spheres denote graph vertices and dark-teal curves denote stored edges. Polynomial annotations belong to the particular displayed graph/diagram and its recorded convention. A source-format label does not certify an automatic source-to-spine correspondence; higher-valence examples do not inherit the unrestricted subcubic interpretation.}
\label{suppfig:input_yamada_examples}
\end{figure}

\subsection{Heterogeneous inputs and the common embedded-graph interface}
\label{supp:input_examples}
The input layer maps each source class to the earliest stage compatible with its mathematical structure. Public adapters include ordered coordinate chains, PDB/PDBx-mmCIF molecular backbones \cite{kg_burley2019rcsb,kg_bourne1997mmcif}, GROMACS \cite{kg_abraham2015gromacs} and LAMMPS \cite{kg_thompson2022lammps} polymer snapshots, and paired node--edge spatial-graph data. Named knots, links, torus types and Artin braid words enter through analytic constructors. Already embedded graphs can enter at the core representation, while accepted diagrams or PD data can bypass three-dimensional geometry and enter the exact-analysis layer directly. The broader input categories in \Figref{fig:input_formats_overview} include biomolecular curves \cite{kg_yeates2007protein_topology}, vascular and neuronal networks \cite{kg_aylward2005spatial_vascular,kg_cuntz2010neuronal_graphs}, polymers \cite{kg_deguchi2017topological_polymers} and engineered networks \cite{kg_shai1999engineering_graphs,kg_peddada2023}. Source labels such as SWC or GraphML describe data that can be externally converted to the common object; they are not a claim that every listed format has a native parser.

Hamiltonian-derived level sets or occupied regions can use application-specific geometry-to-graph routes (\Figpanelref{fig:input_formats_overview}{e}); the companion study introduces a finite-energy Fermi-volume construction \cite{akgun_yan_2026topologicalclassificationknottedgraphs}. Surface meshes enter as geometric objects before a reconstruction is selected. Volumes that are established regular thickenings of embedded networks admit graph spines in the handlebody model (\Figpanelref{fig:input_formats_overview}{f}) \cite{kg_ishii2008handlebody,kg_ishii2012handlebody}. Oriented field curves use the route in panel \textbf{(g)} \cite{kg_helman1989vector_field_topology}, and mathematical spatial graphs can enter directly in panel \textbf{(h)} \cite{kg_kauffman1989spatial,kg_flapan2017spatial}. The interfaces converge on node positions and edge polylines, while the scientific justification for each upstream conversion remains source dependent.

Normalization makes explicit which geometric or topological information is preserved, added or unresolved. Ordered coordinates are validated as finite $(N,3)$ arrays and represented either as an open polyline between endpoint vertices or, for an explicitly closed curve, as a self-loop carrying the ordered path. Closure is not inferred from proximity alone: a direct-closure option adds the final segment explicitly, whereas a metadata-only option records an intention without changing the path. Protein and nucleic-acid adapters select a model, chain and backbone atom trace. Ambiguous multichain inputs require a chain choice, and non-primary alternate locations are skipped rather than merged. Polymer adapters preserve the ordering encoded in the simulation record; LAMMPS snapshots can be selected by molecule identifier and GROMACS coordinates use the stated unit conversion. Paired node--edge CSV inputs preserve attributes, require finite node coordinates and check that supplied edge polylines terminate at their incident vertices.

Surface meshes are not assumed to be graph spines. The surface adapter loads supported meshes into a common \texttt{PyVista PolyData} object, optionally cleans and triangulates the mesh, and checks empty, non-finite or open-boundary geometry before reconstruction. Parsing and representation conversion are therefore separate from the decision that a surface or volume admits an appropriate graph reduction. A closed triangulated surface alone is not evidence of handlebody topology or of a correct ambient embedding of an extracted graph.

After these source-specific choices, edge coordinates, connected components and optional orientation remain accessible independently of the polynomial backend. Biomolecular curves retain ordered conformations; branching networks retain supplied branch incidence; polymers retain their specified curves and links; and engineered networks retain incidence and spatial coordinates. When the spatial embedding itself is the target, the fixed graph can be projected and passed to exact Yamada evaluation, subject to the graph/diagram scope described below.

\begin{figure}[!t]
\centering
\includegraphics[width=\linewidth,keepaspectratio]{GeneralVersion/Figures/InputSkeletonizationBeyondYamada.pdf}
\caption{\textbf{Normalized graph objects support measurements beyond invariant evaluation.}
Examples are \textbf{(a--c)} crambin, ubiquitin and hemoglobin from the RCSB Protein Data Bank \cite{kg_burley2019rcsb}; \textbf{(d)} B-DNA; \textbf{(e)} tRNA; \textbf{(f)} a coiled path; \textbf{(g)} a plain-text cable path; \textbf{(h)} a vascular branch \cite{kg_aylward2005spatial_vascular}; \textbf{(i)} integrated flow paths; \textbf{(j,k)} gyroid and Schwarz-P candidate skeletons; and \textbf{(l)} a Hamiltonian-derived graph, related to the band-geometric input class in Refs.~\cite{akgun_yan_2026topologicalclassificationknottedgraphs,kg_fang2016nodal_line}. Arrows indicate orientation and pale geometry provides context without adding incidence. The word ``Periodic'' in the unchanged label of panel \textbf{(l)} refers to the defining field, not to a verified periodic face-identification operation in the reconstruction. Field-derived graphs illustrate numerical representations rather than certified retractions of every continuous source.}
\label{suppfig:input_skeletonization_beyond_yamada}
\end{figure}

The panels in \SuppFigref{suppfig:input_skeletonization_beyond_yamada} reach the common representation through different routes. Protein and nucleic-acid backbones in \textbf{(a--e)} are ordered traces and enter directly rather than being voxel-skeletonized. The protein records are obtained from the RCSB Protein Data Bank \cite{kg_burley2019rcsb}. The lightweight PDB/mmCIF adapters parse selected backbone records, while Biopython supplies structural-bioinformatics machinery used by the protein-derived layout examples and the broader ecosystem \cite{kg_hamelryck2003biopdb,kg_cock2009biopython}. The explicit paths in \textbf{(f,g)} use the coordinate-chain contract. The vascular branch in \textbf{(h)} supplies node--edge geometry, and the flow paths in \textbf{(i)} retain attached orientation. By contrast, \textbf{(j--l)} illustrate surface/volume-to-graph processing of implicit fields. This distinction prevents direct normalization from being conflated with reconstruction. \SuppFigref{suppfig:input_yamada_examples} complements these examples with particular graph/diagram polynomial outputs.

\FloatBarrier
\subsection{Finite-thickness objects as spatial-graph spines}
\label{supp:handlebody_basis}
Scientific structures may be supplied as surfaces or volumes rather than one-dimensional graphs, as in tubular vascular geometry \cite{kg_aylward2005spatial_vascular} and finite-energy Fermi regions \cite{akgun_yan_2026topologicalclassificationknottedgraphs}. Related network descriptions occur in neuronal morphologies \cite{kg_cuntz2010neuronal_graphs}, engineered systems \cite{kg_shai1999engineering_graphs,kg_peddada2023} and polymeric graph structures \cite{kg_deguchi2017topological_polymers}. When a region is a regular thickening of an embedded network, a spine gives a graph description of its handle structure while discarding local thickness. The computational task is to recover an appropriate representative from sampled data, not to assume that every supplied volume already satisfies this model.

\begin{quote}
\textbf{Abstract graph.} Vertices, edges and their incidence relations, independent of an embedding.
\medskip
\noindent\textbf{Spatial graph.} An embedding of an abstract graph in three-dimensional space, $G\hookrightarrow\mathbb R^3$, for which the embedding is part of the object \cite{kg_kauffman1989spatial,kg_flapan2017spatial}.
\medskip
\noindent\textbf{Terminology.} We use \emph{spatial graph} formally and \emph{knotted graph} to emphasize embedding-sensitive topology, including topologically trivial embeddings.
\end{quote}

The same abstract graph can admit different spatial embeddings. In the structured families of \Figpanelref{fig:knottedgraph_topoly_scaling}{a--f}, the family-specific construction determines incidence while its spatial realization generates a crossing-rich diagram. Some family parameters can also change the size of the underlying graph; crossing growth is not interpreted as proof that every compared diagram has an identical abstract graph. Within a prescribed graph, knotting and linking can change without changing incidence.

Two spatial graphs are \emph{ambient-isotopic} if a continuous deformation of the ambient three-dimensional space carries one to the other without cutting edges or passing nonincident strands through one another. The graph object and the diagrammatic invariant described in \SuppNoteref{supp:yamada_engine} address this embedding information, with the stated valence and vertex-structure conventions.

Finite-thickness topology is related to graphs through the following notions \cite{kg_ishii2008handlebody,kg_ishii2012handlebody}:
\begin{quote}
\textbf{Regular neighbourhood.} A regular neighbourhood $N(G)$ of an embedded graph is a sufficiently small three-dimensional thickening, with vertices expanded to junction regions and edges to tubes.
\medskip
\noindent\textbf{Handlebody.} For connected $G$, $N(G)$ is a compact connected orientable three-manifold obtained from a ball by attaching $1$-handles.
\medskip
\noindent\textbf{Handlebody-link.} Regular neighborhoods of disconnected graph components form a disjoint collection of embedded handlebodies.
\medskip
\noindent\textbf{Spine.} An embedded graph $G\subset\mathcal H$ is a spine when the handlebody $\mathcal H$ deformation retracts onto $G$.
\end{quote}

A handlebody generally has many spines. The relevant equivalence of their regular neighborhoods is described by the following theorem:
\begin{quote}
\textbf{Theorem 2.3 \cite[Theorem~2.3]{kg_ishii2012handlebody}.}
Two spatial graphs represent an equivalent handlebody-link if and only if they are related by a finite sequence of contraction moves and ambient isotopies.
\end{quote}
This theorem concerns spatial moves, not contractions of a graph after its embedding data have been discarded. It neither provides a numerical extraction algorithm nor makes an arbitrary straight-line endpoint update an admissible spatial contraction. In addition, an invariant of a selected graph induces an invariant of its handlebody only if it is compatible with the required changes of spine. Ordinary normalized Yamada values are not assumed to satisfy that additional invariance. A contraction-stable abstract isomorphism class and a spatial-graph invariant therefore answer different questions.

\texttt{KnottedGraph} approaches the reconstruction task through voxelization, thinning, graph extraction and scale-controlled cleanup (\Figpanelref{fig:functionality_overview}{a}; \SuppFigref{suppfig:skeletonization_steps}). The output retains explicit incidence and edge paths for measurement and projection. The inverse benchmark in \SuppSubsecref{supp:handlebody_validation} compares outputs with known generating graphs. Such tests provide empirical checks, while a proved correspondence for an arbitrary analytic source additionally requires discretization and regular-neighborhood evidence.

The same representation principle also points beyond the ordinary handlebody setting. Finite-thickness geometries with distinguished inner boundaries, enclosed voids or nested boundary components motivate a prospective compression-body formulation. A compression body generalizes a handlebody by allowing a distinguished positive boundary together with additional negative boundary components \cite{kg_scharlemann2002heegaard}. This suggests that boundary-resolved collections of spatial graphs and their associated invariants could provide a systematic computational description of more general multi-boundary geometries.

Developing the corresponding mathematical equivalence relations and boundary-resolved invariant sets remains an open theoretical direction. The computational problem is also genuinely more difficult than the handlebody reconstruction treated here. In the handlebody case, the finite-thickness region can be reduced to a purely one-dimensional graph spine, allowing the pipeline to proceed through volumetric thinning, junction reconstruction and embedded-graph cleanup. For geometries with multiple distinguished boundary components, however, a single graph need not contain the full topological information: the relevant deformation retract can involve boundary surfaces together with graph-like attachments. A computational representation must therefore preserve not only graph incidence and spatial embedding, but also which structures attach to which boundary components, their nesting and their relative topology. Standard skeleton-to-graph reduction can erase precisely this information, so new boundary-aware reconstruction, certification and invariant-evaluation procedures are required. The methods introduced here provide a practical starting point for this broader computational theory of finite-thickness topology.

\subsection{Skeleton extraction and topology-aware graph reconstruction}
\label{supp:extraction}
\begin{figure}[!t]
\centering
\includegraphics[width=\linewidth]{GeneralVersion/Figures/SkeletonizationSteps.pdf}
\caption{\textbf{Stages of numerical finite-thickness reconstruction before planar projection.}
\textbf{(a)} Input surface; \textbf{(b)} filled mask; \textbf{(c)} Lee-thinned voxel skeleton; \textbf{(d)} graph extraction using junction-zone candidates and stored edge paths; \textbf{(e)} optional terminal-branch removal; \textbf{(f)} degree-two chain collapse retaining the full polyline; \textbf{(g)} user-controlled short-edge contraction; \textbf{(h)} polyline simplification; and \textbf{(i)} planar projection. The panels isolate cleanup operations. The historical thick-handlebody benchmark used contraction, leaf removal, degree-two simplification and Ramer--Douglas--Peucker smoothing in that order. Chain concatenation removes subdivision without replacing the path; the other geometric updates require separate admissibility checks. Their depiction here is not an unconditional ambient-isotopy certificate.}
\label{suppfig:skeletonization_steps}
\end{figure}

The handlebody model \cite{kg_ishii2012handlebody} supplies a target for the reconstruction in \Figpanelref{fig:functionality_overview}{a}. The numerical task is to represent the graph-like organization of a sampled object while retaining enough geometry for later diagrammatic analysis. The stages in \SuppFigref{suppfig:skeletonization_steps} separate the digital reduction, graph interpretation and optional changes to geometry.

\paragraph{Surface to filled volume: panels (a,b)}
Surface inputs are converted to the intended occupied mask; field- or volume-derived inputs can enter directly at this stage. Surface handling uses \texttt{PyVista} \cite{kg_sullivan2019pyvista}, and level-set surface extraction uses Lewiner marching cubes through \texttt{scikit-image} \cite{kg_lewiner2003marching,kg_vanderwalt2014skimage}. The mask is the digital object on which thinning acts. Its topology depends on resolution, sampling, boundary conventions and the foreground/background adjacency; surface closure alone does not establish a topology-preserving discretization of the continuous source.

\paragraph{Filled volume to voxel skeleton: panel (c)}
Empty margins are cropped to restrict processing to the occupied bounding box. Three-dimensional Lee thinning through \texttt{scikit-image} \cite{kg_lee1994thinning,kg_vanderwalt2014skimage} produces the digital backbone in panel \textbf{(c)}, which is restored to the original coordinate frame. The topological guarantees of a digital thinning rule apply to its specified discrete conventions. They are not a proof that the sampled object is ambient-isotopic to the analytic input, nor that every digital skeleton admits an appropriate one-dimensional graph representation.

\paragraph{Voxel skeleton to raw embedded graph: panel (d)}
The digital skeleton still requires an incidence interpretation. Following voxel-skeleton graph constructions \cite{kg_reinders2000skeletongraph}, voxels touching through a face, edge or corner are first joined using $26$-neighbour adjacency. A diagonal lattice connection is suppressed when the same local connection is represented through an occupied intermediate voxel by two strictly shorter steps. This removes a class of redundant local shortcuts; it is not, by itself, a global embedding certificate.

Degree-two voxels are interpreted as strand interiors. Connected regions of voxels whose degree differs from two identify endpoints or junction candidates, and intervening chains are traced as graph edges. Their ordered three-dimensional coordinates are retained. Entirely degree-two components are represented as loops rather than discarded. Components consisting of isolated voxels or short tree-like pieces must be retained or explicitly classified as outside the intended model, not silently removed merely because they are small. Foreground components and graph components are compared using a stated common convention.

A physical junction can thin to several nearby branch voxels. Conversely, distinct junctions can be spatially close. To address this ambiguity, candidate graphizations are produced at neighboring junction-zone scales. Starting from a zero-radius trace, junction regions are expanded up to a chosen number of voxel hops. Candidates must have nondegenerate edge geometry and avoid anomalously short connections relative to the local edge scale. An expected maximum valence can be supplied as a prior; none is inferred when absent.

Candidate incidence structures are compared by multigraph isomorphism through NetworkX \cite{kg_hagberg2008networkx,kg_cordella2004vf2}. Two consecutive valid isomorphic candidates give the strongest scale-persistence signal, and the smaller-scale representative is retained. Otherwise, the most recurrent valid isomorphism class is selected if it appears at least twice, again choosing its smallest-scale representative. If no class recurs, the direct zero-radius trace is retained rather than promoting an unsupported expanded-scale reconstruction. This is persistence of the extracted abstract graph across a finite set of scales, not persistent homology or a proof of spatial equivalence.

Known valence and component/cycle information can strengthen candidate rejection. The thickened-seed tests use a subcubic prior and junction expansion through four voxel hops. These priors restrict the benchmark problem and are not inferred properties of arbitrary scientific inputs. Repeated incidence and matching Betti numbers remain necessary consistency evidence rather than sufficient evidence of a source-to-spine correspondence.

\paragraph{Leaf removal: panel (e)}
A leaf can be part of the scientific graph and must then be retained. In a justified handlebody-spine setting, contraction of a pendant tree can leave the regular neighborhood equivalent, as allowed by spatial contraction moves \cite{kg_ishii2012handlebody}. This is a change of spine, not preservation of the prescribed abstract graph or of its ordinary Yamada polynomial. A contraction routine must also preserve the relevant connected component: a tree cannot be silently replaced by the empty graph. Biological trees, open polymers and applications in which terminal branches matter do not use universal leaf pruning.

\paragraph{Degree-two chain collapse: panel (f)}
Consecutive degree-two graph segments are concatenated into a single edge while retaining the entire ordered spatial path. This removes redundant subdivision without replacing the embedded curve by a chord. Closed degree-two components remain represented as loops. The distinction between concatenation and geometric simplification is important for the later embedding-sensitive calculation.

\paragraph{Short-edge contraction: panel (g)}
A short edge may indicate voxel-scale fragmentation of a junction, but short length alone does not establish that the implemented motion is a valid spatial contraction. The numerical operation combines endpoint vertices and relinks incident polylines. To infer preservation of the embedded regular neighborhood, the motion must meet the hypotheses of an admissible spatial contraction, including separation from nonincident geometry. An abstract contraction or unchanged cycle rank cannot supply that evidence. The threshold is therefore a user-controlled modeling and cleanup parameter; changing it may change the selected graph and its signature. A bounded search that finds no admissible sequence reports only that equivalence was not established within the search budget.

\paragraph{Geometric smoothing: panel (h)}
Ramer--Douglas--Peucker simplification \cite{kg_ramer1972,kg_douglas1973}, implemented through \texttt{fastrdp}, reduces voxel-scale staircase structure along retained polylines. Endpoints are normalized to graph-vertex coordinates and consecutive duplicate samples are removed. These operations can improve numerical conditioning while preserving incidence. The cited simplification algorithm does not supply an unconditional ambient-isotopy guarantee: a shortcut can meet another strand even when its endpoints and abstract graph are unchanged. Embedding preservation must therefore be checked through admissible swept geometry or other appropriate separation criteria, or left unasserted for the processed result. The independently guarded optional layout route is discussed in \SuppNoteref{supp:repulsive_layout}.

\paragraph{Embedded graph to planar diagram: panel (i)}
The output carries explicit node positions and ordered edge polylines. It can be used for geometric and graph measurements without a polynomial calculation. For embedding-sensitive analysis, a generic projection resolves crossings and produces PD data. This representation change is described in \SuppNoteref{supp:projection}. An exact polynomial of the resulting diagram is not retrospective proof that the earlier volume reduction or geometric cleanup was topology preserving.

The sequence therefore separates digital reduction \textbf{(a--c)}, incidence reconstruction \textbf{(c,d)}, optional graph/geometric processing \textbf{(e--h)} and diagram construction \textbf{(i)}. The inverse benchmark tests selected outputs against known generators. The detailed limitations of each correspondence are collected in \SuppNoteref{supp:limitations}.
